\chapter{Die Bühne und die Bibel}
\quellseite{10}

Die berühmte schwedische Sängerin Jenny Lind, genannt die „schwedische Nachtigall“, hatte großen Erfolg als Opernsängerin, und das Geld strömte ihr nur so zu. Gleichwohl verließ sie die Opernbühne, als sie auf der Höhe ihrer Kunst stand, und kehrte nie zu ihr zurück. Alles gab sie daran, Geld, Ruhm, Beifall der Menge, und zog sich in die Stille zurück.

Eine englische Freundin kam einst zu ihr, als sie am Meeresstrand saß, ihre liebe Lutherische Bibel auf ihrem Schooß hielt und den Untergang der Sonne betrachtete. Beide geriethen ins Gespräch und die Freundin fragte sie, was so viele Bewunderer der Sängerin sich oft im Stillen gefragt hatten: „O Madame Goldschmidt“ (so hieß Jenny Lind nach ihrer Verheirathung), „wie kam es doch, daß Sie die Bühne gerade auf dem Höhepunkt Ihres Erfolges verließen?“

„Was anders konnte ich thun“, lautete die einfache Antwort, „wenn jeder Tag meines Bühnenlebens mich weniger an diese“ (dabei legte sie den Finger auf die Bibel) „und gar nicht an jene“ (dabei deutete sie auf die untergehende Sonne) „denken ließ?“
